
\setcounter{demsode}{5} %%% THỨ TỰ ĐỀ THẦY - CÔ ĐÃ NHẬN

\begin{name}
	{Biên soạn: Thầy Lê Hải Phụng\\
	Phản biện: Thầy Nguyễn Hữu Đức}
	{ÔN TẬP HỌC KỲ II KHỐI 10 NĂM 2024}
\end{name}

\cautn
\Opensolutionfile{ans}[ans/0-CTST-TN-SO-5]
\setcounter{ex}{0}

%%% CÂU 1
%%%============EX_1==============%%%
\begin{ex}%[0D8N1-1]
	Từ thành phố $A$ đến thành phố $B$ có 5 cách đi bằng đường bộ, 3 cách đi bằng đường thủy và 2 cách đi bằng đường hàng không. Hỏi có bao nhiêu cách đi từ thành phố $A$ đến thành phố $B$?
	\choice
	{$16$}
	{$15$}
	{\True $10$}
	{$30$}
	\loigiai{}
\end{ex}

%%%============EX_1==============%%%
\begin{ex}%[0D8N2-3]
	Từ các chữ số $1,2,3,4,5,6,7$ lập được bao nhiêu số tự nhiên gồm hai chữ số khác nhau?
	\choice
	{$C_7^2$}
	{$2^7$}
	{$7^2$}
	{\True $A_7^2$}
	\loigiai{}
\end{ex}
%%%============EX_2==============%%%
\begin{ex}%[0H9N4-1]
	Trong mặt phẳng tọa độ $Ox y$, cho đường tròn $(C)\colon (x+1)^2+(y-3)^2=4$. Tọa độ tâm $I$ của đường tròn $(C)$ là
	\choice
	{$I(1;-3)$}
	{\True $I(-1; 3)$}
	{$I(1; 3)$}
	{$I(-1;-3)$}
	\loigiai{}
\end{ex}
%%%============EX_3==============%%%
\begin{ex}%[0H9N3-2]
	Cho phương trình đường thẳng $\Delta: 3x+4y-5=0$. Tìm một véctơ pháp tuyến của đường thẳng $\Delta$.
	\choice
	{$\vec{n}=(4;-3)$}
	{$\vec{n}=(4; 3)$}
	{$\vec{n}=(-4; 3)$}
	{\True $\vec{n}=(3; 4)$}
	\loigiai{}
\end{ex}
%%%============EX_4==============%%%
\begin{ex}%[0D0H2-5]
	Một hộp có 4 bi đỏ, 3 bi xanh, 2 bi vàng. Lấy ngẫu nhiên 3 bi. Tính xác suất để 3 bi lấy ra có ít nhất một bi đỏ.
	\choice
	{$\dfrac{2}{7}$}
	{\True $\dfrac{37}{42}$}
	{$\dfrac{3}{4}$}
	{$\dfrac{10}{21}$}
	\loigiai{}
\end{ex}
%%%============EX_5==============%%%
\begin{ex}%[0H9N5-2]
	Trong mặt phẳng $Ox y$, phương trình nào sau đây là phương trình chính tắc của một Elip?
	\choice
	{$x^2+y^2=1$}
	{\True $\dfrac{x^2}{2^2}+\dfrac{y^2}{1^2}=1$}
	{$\dfrac{x^2}{2}+\dfrac{y^2}{2}=1$}
	{$\dfrac{x^2}{4}+\dfrac{y^2}{9}=1$}
	\loigiai{}
\end{ex}
%%%============EX_6==============%%%
\begin{ex}%[0D8N1-2]
	Trong tủ quần áo của bạn An có 4 chiếc áo khác nhau và 3 chiếc quần khác nhau. Hỏi bạn An có bao nhiêu cách để chọn 1 bộ quần áo để mặc?
	\choice
	{$27$}
	{$64$}
	{\True $12$}
	{$7$}
	\loigiai{}
\end{ex}
%%%============EX_7==============%%%
\begin{ex}%[0H9N1-1]
	Điểm nào sau đây nằm trên đường thẳng $d:\heva{&x=1+2t \\
		&y=2-t} \quad(t \in \mathbb{R})$?
	\choice
	{$P(-1;-2)$}
	{$Q(5; 1)$}
	{\True $M(-1; 3)$}
	{$N(2;-1)$}
	\loigiai{}
\end{ex}
%%%============EX_8==============%%%
\begin{ex}%[0D7N1-2]
	Số giá trị nguyên âm của $x$ để tam thức tam thức bậc hai $f(x)=2x^2+7x-9$ nhận giá trị âm là
	\choice
	{$3$}
	{$5$}
	{$6$}
	{\True $4$}
	\loigiai{}
\end{ex}
%%%============EX_9==============%%%
\begin{ex}%[0D7N3-2]
	Phương trình $\sqrt{-x^2+4x}=2x-2$ có bao nhiêu nghiệm?
	\choice
	{$3$}
	{$0$}
	{$2$}
	{\True $1$}
	\loigiai{}
\end{ex}
%%%============EX_10==============%%%
\begin{ex}%[0H9N3-1]
	Trong mặt phẳng $Ox y$, đường thẳng đi qua điểm $A(2;-4)$ và nhận $\vec{u}=(-4; 3)$ là vec-tơ chỉ phương có phương trình tham số là
	\choice
	{\True $\heva{&x=2-4t \\
			&y=-4+3t
			&}$}
	{$\dfrac{x-2}{-4}=\dfrac{y+4}{3}$}
	{$\heva{&x=-4+2t \\
			&y=3-4t
			&}$}
	{$\heva{&x=-2-4t \\
			&y=4+3t
			&}$}
	\loigiai{}
\end{ex}
%%%============EX_11==============%%%
\begin{ex}%[0D0N1-3]
	Rút ngẫu nhiên cùng lúc 3 quân bài từ cỗ bài lơ khơ 52 quân, số phần tử của không gian mẫu $n(\Omega)$ bằng
	\choice
	{\True $22\, 100$}
	{$156$}
	{$132\, 600$}
	{$140\, 608$}
	\loigiai{}
\end{ex}

\Closesolutionfile{ans}
\indapan{6}{ans/0-CTST-TN-SO-5}

\cauds
\Opensolutionfile{ans}[ans/0-CTST-DS-SO-5]
\setcounter{ex}{0}
Câu trắc nghiệm đúng sai. Học sinh trả lời từ câu 1 đến câu 4. Trong mỗi ý a), b), c), d) ở mỗi câu, học sinh chọn đúng hoặc sai.
%%%============EX_1==============%%%
\begin{ex}%[0D8H1-3]
	Hình sau đây biểu diễn các con đường một chiều nối các thành phố $A, B$ và $C$, khi đó:
\begin{center}
		\begin{tikzpicture}[every node/.style={circle, fill=red!30},scale=1]
		\tikzset{decoration={markings,mark=between positions 0.5 and 0.5 step 1 with {\draw[-stealth] (0,0)--++(0:.1);}}}
		\path (0,0) node(A){A} (2.3,0) node(B){B} (4,0) node(C){C};
		\draw[postaction={decorate}] (A)--(B);
		\draw[postaction={decorate},out=150,in=30] (B) to (A);
		\draw[postaction={decorate},out=-150,in=-30] (B) to (A);
		\draw[postaction={decorate}] (C)--(B);
		\draw[postaction={decorate},out=30,in=150] (B) to (C);
		\draw[postaction={decorate},out=-150,in=-30] (C) to (B);
		\draw[postaction={decorate}] (A.50) .. controls (1,1) and (3,1) .. (C.130);
		\draw[postaction={decorate}] (A.-50) .. controls (1,-1) and (3,-1) .. (C.-130);
	\end{tikzpicture}
\end{center}
	\choiceTF
	{\True Có 2 cách di chuyển từ thành phố $A$ đến thành phố $C$ mà không đi qua thành phố B}
	{\True Có 1 cách di chuyển từ thành phố $A$ đến thành phố $C$ mà đi qua thành phố B}
	{Có 3 cách đi từ thành phố $A$ đến thành phố $B$ mà không đi qua thành phố C}
	{Có 3 cách đi từ thành phố $A$ đến thành phố $C$ rồi quay trở lại thành phố A}
	\loigiai{}
\end{ex}
%%%============EX_2==============%%%
\begin{ex}%[DA-HK2-K10,11-2024]
	Cho biểu thức $f(x)=(3x-1)\left(3x^2-4x+1\right)$. Khi đó
	\choiceTF
	{\True $f(x)=0\Leftrightarrow\hoac{&x=-\dfrac{1}{3} \\&x=1&}$}
	{\True Với $x \in\left(-\infty; \dfrac{1}{3}\right) \cup\left(\dfrac{1}{3}; 1\right)$ thì $f(x) < 0$}
	{Với $x \in(1;+\infty)$ thì $f(x) < 0$}
	{\True Bảng xét dấu của biểu thức là \\
		\begin{tikzpicture}
			\tkzTabInit[lgt=3,espcl=2]
			{$x$/0.7,$3x-1$/0.7,$3x^2-4x+1$/0.7,$f(x)$/.7}
			{$-\infty$,$\tfrac{1}{3}$,$1$,$+\infty$}
			\tkzTabLine{,-,t,+,z,+,}
			\tkzTabLine{,+,z,-,z,+,}
			\tkzTabLine{,-,z,-,z,+,}
		\end{tikzpicture}
	}
	\loigiai{
		
	}
\end{ex}
%%%============EX_3==============%%%
\begin{ex}%[DA-HK2-K10,11-2024]
	Cho biểu thức $f(x)=\dfrac{1}{x^2-2x-12}$. Khi đó:
	\choiceTF
	{$f(x)=0\Leftrightarrow x=1+\sqrt{13}$ hoặc $x=1-\sqrt{13}$}
	{Với $x \in(1-\sqrt{13}; 1+\sqrt{13})$ thì $f(x) > 0$}
	{Với $x \in(-\infty; 1-\sqrt{13}) \cup(1-\sqrt{13};+\infty)$ thì $f(x) < 0$}
	{\True Bảng xét dấu của biểu thức là\\
		\begin{tikzpicture}
			\tkzTabInit[lgt=3,espcl=2]
			{$x$/0.7,$x^2-2x-12$/0.7,$f(x)$/.7}
			{$-\infty$,$1-\sqrt{13}$,$1-\sqrt{13}$,$+\infty$}
			\tkzTabLine{,+,z,-,z,+,}
			\tkzTabLine{,+,d,-,d,+,}
	\end{tikzpicture}}
	\loigiai{
		
	}
\end{ex}
%%%============EX_4==============%%%
\begin{ex}%[0H9H3-2]
Xét tính đúng, sai của các khẳng định sau:
\choiceTF
{$\Delta$ qua $M(2;-3)$ và vuông góc với $AB$ và $A(1,5), B(-4,7)$, khi đó phương trình tổng quát của $\Delta$ là: $-5x+2y+16=0$}
{$\Delta$ đi qua $A(-1,2)$ và $B(3,-1)$, khi đó phương trình tổng quát của $\Delta$ là: $3x+4y-5=0$}
{$\Delta$ qua $A(-3,5)$ và $\Delta \perp d: x-2y+3=0$, khi đó phương trình tổng quát của $\Delta$ là: $x+y-2=0$}
{$\Delta$ qua $A(-1,2)$ và $\Delta \parallel d: x=3$, khi đó phương trình tổng quát của $\Delta$ là: $x+y-1=0$}
\loigiai{}
\end{ex}

\Closesolutionfile{ans}
\indapan{4}{ans/0-CTST-DS-SO-5}

\caukq
\Opensolutionfile{ans}[ans/0-CTST-KQ-SO-5]
\setcounter{ex}{0}
 Trắc nghiệm lựa chọn câu trả lời ngắn. Học sinh trả lời từ câu 1 đến câu 6.
%%% CÂU 1
\begin{ex}%[DA-HK2-K10,11-2024 - TÊN THẦY CÔ ???] %[ID6]
	Có bao nhiêu giá trị nguyên của $m$ để phương trình $x^{2}-(m+1) x+3 m-5=0$ vô nghiệm?	
	\shortans[]{$3$}
	\loigiai{
	}	
\end{ex}

%%% CÂU 2
\begin{ex}%[DA-HK2-K10,11-2024 - TÊN THẦY CÔ ???]%[0D8V2-4]
	Ở một Đoàn trường phổ thông có $5$ thầy giáo, $4$ cô giáo và $8$ học sinh. Có $X$ cách chọn ra một đoàn công tác gồm $7$ người trong đó có $1$ trưởng đoàn là thầy giáo, $1$ phó đoàn là cô giáo và đoàn công tác có ít nhất $4$ học sinh. Tính tổng các chữ số của $X$.
	\shortans[]{$12$}	
	\loigiai{
		
	}
\end{ex}

%%% CÂU 3
\begin{ex}%[DA-HK2-K10,11-2024 - TÊN THẦY CÔ ???]%[0D8V3-4]
	Tìm hệ số của $x^{4}$ trong khai triển $(2-3x)^{5}+2(3x+1)^{6}$.
	\shortans{$3240$}
	\loigiai{
		
	}
\end{ex}

%%% CÂU 4
\begin{ex}%[DA-HK2-K10,11-2024 - TÊN THẦY CÔ ???] %[0D7H1-5]
	Một người muốn uốn tấm tôn phẳng hình chữ nhật có bề ngang $32 \mathrm{~cm}$, thành một rãnh dẫn nước bằng cách chia tấm tôn đố thành ba phần rồi gấp hai bên lại theo một góc vuông như hình vẽ. Biết rằng diện tích mặt cắt ngang của rãnh nước phải lớn hơn hoặc bằng tổng $120 \mathrm{~cm}^{2}$. Hỏi độ cao tối thiểu của rãnh dẫn nước là bao nhiêu cm?
	\begin{center}
		\begin{tikzpicture}[>=stealth,font=\footnotesize,scale=.7]
			\path (0,0) coordinate (A) ++(0:1.5) coordinate (B) ++(0:2)  coordinate (C) ++(0:1.5) coordinate (D)
			(60:5) coordinate (A1) ++(0:1.5) coordinate (B1) ++(0:2)  coordinate (C1) ++(0:1.5) coordinate (D1);
			\draw (A)--(D)--(D1)--node[above]{$x$}(C1)--(B1)--node[above]{$x$}(A1)--cycle (B)--(B1) (C)--(C1);
			\draw[<->,cyan] (A1)++(90:.5)--node[above,black]{$32$ cm}++(0:5);
			\foreach \i in {A,B,C,D,A1,B1,C1,D1} \fill (\i) circle (1.2pt);
			\begin{scope}[xshift=8cm]
				\path (0,0) coordinate (A) ++(0:2) coordinate (B) ++(90:1.5)  coordinate (C) ++(180:2) coordinate (D)
				(60:2.5) coordinate (A1) ++(0:2) coordinate (B1) ++(90:1.5)  coordinate (C1) ++(180:2) coordinate (D1)
				(60:5) coordinate (A2) ++(0:2) coordinate (B2) ++(90:1.5)  coordinate (C2) ++(180:2) coordinate (D2);
				\fill[gray,opacity=.4] (A1)--(B1)--(C1)--(D1)--cycle;
				\draw (A)--(B)--(C)--(D)--node[left]{$x$}cycle (D)--(D2)--(C2)--(C) (B)--(B2)--node[right]{$x$}(C2) (B1)--(C1)--(D1);
				\draw[dashed] (A)--(A2)--(B2) (D1)--node[left,pos=.7]{$x$}(A1)--(B1) (A2)--(D2);
				\draw[<-] (barycentric cs:D1=1,C1=2)--++(135:2.5) node[above]{mặt cắt};
				\foreach \i in {A,B,C,D,A1,B1,C1,D1,A2,B2,C2,D2} \fill (\i) circle (1.2pt);
			\end{scope}
		\end{tikzpicture}
	\end{center}
	\shortans[]{$6$}
	\loigiai{
	}	
\end{ex}

%%% CÂU 5
\begin{ex}%[DA-HK2-K10,11-2024 - TÊN THẦY CÔ ???]%[0H9H4-2]
	Viết phương trình đường tròn $(C)\colon x^2+y^2-2ax-2by+c=0$ biết $(C)$ đi qua hai điểm $A(3;3)$, $B(7;-1)$  và có tâm nằm trên đường thẳng $d\colon x+y-2=0$. Khi đó $a+b+c$ bằng bao nhiêu?
	\shortans[]{$-4$}
	\loigiai{
	
	}
\end{ex}

%%% CÂU 6
\begin{ex}%[DA-HK2-K10,11-2024 - TÊN THẦY CÔ ???]%[0D0V2-5]
	Một hộp đựng $15$ viên bi, trong đó có $7$ viên bi xanh và $8$ viên bi đỏ. Lấy ngẫu nhiên ba viên bi (không kể thứ tự ra khỏi hộp). Tính xác suất để trong ba viên bi lấy ra có ít nhất một viên bi đỏ (làm tròn kết quả đến chữ số thập phân thứ hai).
	\shortans[]{$0{,}92$}
	\loigiai{
		Số phần tử của không gian mẫu $n(\Omega)=\mathrm{C}_{15}^3=455$.\\
		Gọi $A$ là biến cố \lq\lq trong ba viên bi lấy ra có ít nhất một viên bi đỏ \rq\rq.\\
		Suy ra $\overline{A}$ là biến cố \lq\lq trong ba viên bi lấy ra không có bi đỏ nào\rq\rq.\\
		Chọn $3$ bi xanh trong $7$ bi xanh có $\mathrm{C}_7^3=35$ cách.\\
		$$\Rightarrow n(\overline{A})=35\Rightarrow \mathrm{P}(\overline{A})=\dfrac{n(\overline{A})}{n(\Omega)}=\dfrac{1}{13}.$$
		Vậy $\mathrm{P}(A)=1-\mathrm{P}(\overline{A})=\dfrac{12}{13} \approx 0{,}92$.
	}
\end{ex}
\Closesolutionfile{ans}
\indapan{6}{ans/0-CTST-KQ-SO-5}